%% ===================================================================
\section{Introduction}
\label{sec:intro}
%% ===================================================================

Let $G$ be a finite group and let $\mathcal C\subseteq G$ be a
conjugation-closed set generating $G$, for instance a single conjugacy
class. Label the $2^m$ leaves of a balanced binary tree of depth $m$ by
elements of $\mathcal C$ and let every internal vertex take the commutator
$[x,y]=xyx^{-1}y^{-1}$ of its two children. Expanding the root value as a
word in the leaf labels gives a product of $4^m$ occurrences, each a leaf
label or its inverse. Now read every occurrence independently: with
probability $1/2$ it contributes its element, otherwise the identity. The
output is a random product of $4^m$ independent lazy factors, that is, a
lazy random walk on $G$ whose steps are prescribed by the tree rather than
chosen freely. We ask when every such walk contracts towards the uniform
distribution, and what determines the depth at which this begins.

Two features distinguish the construction from a word map. First,
repeated occurrences of a label, including the inverse copies inside each
commutator, receive independent coin tosses, so the output is not the
value of a word at random arguments; in particular it is not confined to
the derived subgroup, and one-dimensional characters are contracted by
the individual lazy factors. Second, only \emph{valid} configurations are
considered, those whose deterministic commutator value at the root is
non-trivial. This constraint is what brings the derived series into the
problem: a configuration whose leaves all lie in a solvable subgroup $H$
has trivial root value as soon as the depth reaches the derived length of
$H$.

We call a labelled tree a \emph{configuration}, and say that it
\emph{contracts strictly} if its Fourier operator has norm less than one
in every non-trivial irreducible representation of $G$, that is, if the
walk has a spectral gap on the orthogonal complement of the constants. A
subgroup of $G$ is \emph{$\mathcal C$-generated} if it is generated by
its intersection with $\mathcal C$. Precise definitions are given in
\Cref{sec:prelim}.

\subsection*{The solvability dichotomy}
The answer is governed by solvability, and the depth by derived length.

\begin{lettertheorem}[Solvability dichotomy and exact level]
\label{ithm:A}
Let $G$ be a finite non-solvable group and $\mathcal C$ a
conjugation-closed generating set of $G$.
\begin{enumerate}[nosep,leftmargin=1.9em,label=\textup{(\arabic*)}]
\item If some proper $\mathcal C$-generated subgroup of $G$ is
non-solvable, then at every depth there is a valid configuration that
fails to contract in some non-trivial irreducible representation.
\item If every proper $\mathcal C$-generated subgroup of $G$ is solvable,
let $D$ be the largest derived length of such a subgroup. Then every
valid configuration of depth at least $D$ contracts strictly, and at
every depth less than $D$ some valid configuration does not.
\end{enumerate}
\end{lettertheorem}

Thus $(G,\mathcal C)$ has a \emph{uniform strict-contraction level}, a
depth from which every valid configuration contracts strictly, if and
only if every proper $\mathcal C$-generated subgroup is solvable, and
the level is then exactly $D$ (\Cref{thm:dichotomy}). Equivalently, $D$
is the exact depth from which a non-trivial iterated commutator certifies
that its entries generate $G$. Transpositions illustrate both cases. In
$S_5$ every proper subgroup generated by transpositions is solvable, the
largest derived length is that of $S_4$, and the level is $3$. In $S_n$
with $n\ge6$ the point stabiliser $S_{n-1}$ is a proper non-solvable
subgroup generated by transpositions, and at every depth it supplies
valid configurations supported in a proper subgroup (\Cref{cor:Sn}).
Only $\mathcal C$-generated subgroups matter: for the involutions of
$A_5$ the level is $2$, because the involutions of $A_4$ generate only
$V_4$ (\Cref{cor:A5}). \Cref{tab:synth} lists further levels.

Applied to all conjugacy classes at once, \Cref{ithm:A} characterises
the minimal simple groups classified by Thompson~\cite{Thompson1968}.

\begin{corollary}[Minimal simple groups]
\label{icor:minsimple}
A finite non-abelian simple group $G$ is minimal simple, that is, all of
its proper subgroups are solvable, if and only if every non-trivial
conjugacy class of $G$ has a uniform strict-contraction level. In that
case every class has level at most $5$, and the involution class of
$\PSL_3(3)$ has level exactly $5$
\textup{(\Cref{cor:minsimple,prop:psl33-level})}.
\end{corollary}

\begin{table}[t]
\centering\small
\begin{tabular}{p{.29\textwidth}p{.40\textwidth}p{.20\textwidth}}
\toprule
Group and leaf set $\mathcal C$ & Decisive proper $\mathcal C$-generated subgroups & Level \\
\midrule
$S_5$, transpositions & $S_4$, of derived length $3$ & $3$ \\
$S_n$, transpositions, $n\ge6$ & $S_{n-1}$, non-solvable & none \\
$A_5$, involutions & $S_3$ and $D_{10}$, of derived length $2$ & $2$ \\
$\PSL_2(7)$, involutions & $S_4$ & $3$ \\
$\PSL_3(3)$, involutions & $3^2{:}\GL_2(3)$, of derived length $5$ & $5$ \\
$\SL_3(3)$, transvections & $3^2{:}\SL_2(3)$, of derived length $4$ & $4$ \\
$\PGL_2(q)$, all involutions & $\PSL_2(q)$, non-solvable & none \\
$\PGL_2(q)$, exterior involutions & subfield subgroups of odd degree; otherwise affine, dihedral or $S_4$ & $2$, $3$ or none (\Cref{ithm:C}) \\
\bottomrule
\end{tabular}
\smallskip
\caption{Uniform strict-contraction levels for some leaf sets; $q\ge5$ is odd
(\Cref{sec:synthesis,sec:rankone}).}
\label{tab:synth}
\end{table}

The proof of \Cref{ithm:A} has two ingredients. The first is elementary.
The Fourier operator of a configuration is a product of lazy averages
$\tfrac12(I+\rho(g))$, and such a product has norm one exactly when the
leaves have a common non-zero fixed vector; hence a configuration
contracts strictly if and only if its leaves generate $G$
(\Cref{thm:criterion}). The second is the normal-subset word-value lemma
of Fern\'andez-Alcober and
Pintonello~\cite[Lemma 2.1(ii)]{FernandezPintonello2023}: for a
conjugation-closed generating subset of a group $H$, the depth-$k$
commutator values of its elements generate $H^{(k)}$. Together they show
that validity forces the leaves out of every proper solvable
$\mathcal C$-generated subgroup once the depth reaches its derived
length, and never forces them out of a non-solvable one. Contraction has
to be tested in every irreducible representation: a proper subgroup that
is transitive in some permutation action is invisible to the
corresponding permutation representation (\Cref{rem:std}).

\Cref{ithm:A} is related to two lines of work. Probabilistic generation
asks how often a random tuple escapes every proper subgroup: for finite
simple groups almost always, by Dixon, Kantor--Lubotzky and
Liebeck--Shalev~\cite{Dixon1969,KantorLubotzky1990,LiebeckShalev1995},
with refinements for elements of prescribed
type~\cite{LiebeckShalev1996,Burness2019} and beyond simple
groups~\cite{HarperQuick2026}. Here the escape is forced by validity
rather than by chance, and the derived length records when. Secondly,
solvability is detected through subgroups generated by few or prescribed
elements in Thompson's classification~\cite{Thompson1968}, in the
descriptions of the solvable radical by two-generated subgroups and by
conjugates~\cite{GKPS2006,GGKP2008}, in two-variable
laws~\cite{Bandman2006,BrayWilsonWilson2005} and in probabilistic
identities~\cite{LarsenShalev2016}. \Cref{ithm:A} is a statement of this
kind with an analytic conclusion: solvability of the proper
$\mathcal C$-generated subgroups is equivalent to the eventual onset of a
uniform spectral gap, and their derived lengths give the exact onset.

\subsection*{A sharp angle between octahedral subgroups}
When no common fixed vector exists, contraction is governed by angles
between fixed spaces. For subgroups $A,B\le G$ and a representation
$\rho$, let $P_A^\rho,P_B^\rho$ be the orthogonal projections onto the
fixed spaces; then $\opn{P_A^\rho P_B^\rho}$ is the cosine of the angle
between these spaces, and for conjugate subgroups it is the norm of a
double coset in a Hecke algebra (\Cref{cor:double-coset-angle}). Such
angles control alternating
projections~\cite{KayalarWeinert1988,BadeaGrivauxMuller2010} and give
Kazhdan constants from arrangements of
subgroups~\cite{Kassabov2011}. Our main concrete result computes one of
them sharply, uniformly in a family of rank one.

For a prime $q\equiv\pm3\pmod 8$ the subgroups of $\PGL_2(q)$ isomorphic
to $S_4$ are not contained in $\PSL_2(q)$; we call them \emph{exterior
octahedral} subgroups. They form a single conjugacy
class~\cite{Huppert1967,King2005,Giudici2007}. Reducing Hamilton's
quaternions modulo $q$ produces an explicit such subgroup $A_q$ and a
conjugate $B_q=hA_qh^{-1}$, where $h=i+j+k$.

\begin{lettertheorem}[The natural octahedral pair]
\label{ithm:B}
Let $q\ne3$ be a prime with $q\equiv\pm3\pmod8$ and put $G_q=\PGL_2(q)$.
Then $A_q\cap B_q\cong S_3$ and $\langle A_q,B_q\rangle=G_q$. Every pair
$A_1,A_2$ of exterior octahedral subgroups of $G_q$ with
$|A_1\cap A_2|=6$ is simultaneously conjugate to $(A_q,B_q)$, and
\[
 \max_{\rho\in\widehat{G_q},\,\rho\ne\mathbf1}
              \opn{P_{A_1}^\rho P_{A_2}^\rho}=c_q\le\frac{\sqrt3}{2},
 \qquad c_q\longrightarrow\frac{\sqrt3}{2}\quad(q\to\infty),
\]
where $c_q$ is the value attained by the pair $(A_q,B_q)$
\textup{(\Cref{prop:octahedral-cap6,thm:octahedral-sharp})}.
\end{lettertheorem}

The mechanism is a compression identity. The four quaternions
$\pm i\pm j\pm k$ of norm $3$ represent the double coset $A_qhA_q$, so
the Hecke operator of the $4$-regular Lubotzky--Phillips--Sarnak
graph~\cite{LPS1988} $X^{3,q}$, restricted to the $A_q$-fixed vectors of
each irreducible representation, \emph{is} the angle operator
$P_{A_q}\rho(h)P_{A_q}$. The Ramanujan bound $2\sqrt3$, divided by the
degree $4$, is $\sqrt3/2=\cos(\pi/6)$, and the Alon--Boppana bound on
the quotient graphs gives optimality. Geometrically, the smallest angle
between the two fixed spaces over all non-trivial irreducible
representations is at least $\pi/6$ and tends to $\pi/6$. Combined with
an estimate inside $S_4$, the angle bound yields an explicit uniform
contraction constant for the corresponding configurations
(\Cref{cor:octahedral-block}).

By \Cref{prop:octahedral-cap6} two distinct exterior octahedral
subgroups meet in a subgroup of order at most $6$, so \Cref{ithm:B}
settles the largest intersection in the following problem, which was
raised in the preliminary
preprint~\cite[Conjecture 6.9]{ZermenoRankOneZenodo}.

\begin{conjecture}[Uniform angle for exterior octahedral pairs]
\label{conj:allpairs}
For primes $q\equiv\pm3\pmod8$, $q\ne3$,
\[
 \sup_{\substack{A,B\cong S_4\ \mathrm{exterior}\\
                 \langle A,B\rangle=\PGL_2(q)}}\ \
 \max_{\rho\ne\mathbf1}\opn{P_A^\rho P_B^\rho}
 \le\frac{\sqrt3}{2}+o(1)\qquad(q\to\infty).
\]
\end{conjecture}

The intersection orders $1,2,3,4$ remain open; numerical all-pairs
maxima for five primes are recorded in \Cref{tab:allpairs-evidence}, and
a qualitative bound $1-b$ with an absolute $b>0$ holds outside a set of
primes of relative density zero (\Cref{cor:almost-all-primes}).

\subsection*{Exact levels in rank one}
The subgroup criterion can be evaluated in infinite families. An
involution of $\PGL_2(q)$ is \emph{exterior} if it lies outside
$\PSL_2(q)$.

\begin{lettertheorem}[Rank-one involutions]
\label{ithm:C}
Let $q\ge5$ be odd. In $\PGL_2(q)$ the set of all involutions has no
uniform strict-contraction level. The exterior class has a level if and
only if $\F_q$ has no proper subfield $\F_{q_0}$ with $q_0\ge5$ and
$[\F_q:\F_{q_0}]$ odd; the level is then $3$ for $q\equiv\pm3\pmod8$
and $2$ for $q\equiv\pm1\pmod8$. The involution class of $\PSL_2(q)$ is classified in
\Cref{thm:rankone-classification}.
\end{lettertheorem}

The level depends on the class and not only on the group: in $\PSL_3(3)$
the involutions have level $5$ and the transvections level $4$. For
transvections in $\SL_n(q)$, $n\ge3$, the level is $3$ for $\SL_3(2)$
and $4$ for $\SL_3(3)$, and in every other case a proper non-solvable
subgroup generated by transvections obstructs
(\Cref{prop:transvection-levels}). For a generating conjugacy class of
$\GL_2(\Z/p^r)$, $p\ge5$, the dichotomy is inherited from the reduction
modulo $p$, and at $r=2$ the level increases by at most one
(\Cref{prop:congruence-class}).

\subsection*{Rates}
For a fixed group, strict contraction gives a positive gap by
finiteness; in a family of groups the gap may tend to zero. Write $U_G$
for the uniform distribution on $G$ and $\mu_\tau$ for the output
distribution of a configuration $\tau$.

\begin{lettertheorem}[Uniform rates and optimal order of depth]
\label{ithm:D}
For the classes with a level in \Cref{ithm:C} and
\Cref{thm:rankone-classification}, and for every non-trivial conjugacy
class of every minimal simple group, there are absolute constants
$c,B>0$ such that every valid configuration $\tau$ of depth $m\ge D$
satisfies
\[
 \dtv(\mu_\tau,U_G)\le\tfrac12\sqrt{|G|-1}
                  \exp\bigl(-c4^m/(\log|G|)^B\bigr).
\]
For a fixed error in $(0,1/2)$, the least depth guaranteeing this error
for all valid configurations has order $\Theta(\log\log|G|)$ as
$|G|\to\infty$
\textup{(\Cref{thm:family-rate,prop:depth-lower,cor:family-examples})}.
\end{lettertheorem}

The upper bound converts the bounded-rank diameter theorem of
Pyber--Szab\'o~\cite{PyberSzabo2016} into a contraction estimate at depth
$D$ and amplifies it through the tree; the lower bound holds because the
output at depth $m$ is supported on at most $2^{4^m}$ elements. For
involution leaves at bounded depth, a uniform contraction gap is
equivalent to a uniform spectral gap for the associated lazy Cayley
walks (\Cref{cor:bounded-depth-gap}), so the quantitative problem is one
of expansion for particular generating sets, in the setting of
quasirandomness, product growth and the Bourgain--Gamburd
method~\cite{Gowers2008,BabaiNikolovPyber2008,Helfgott2008,BGT2011,
BourgainGamburd,BourgainGamburdJEMS2008,BourgainVarju2012,SalehiVarju,
BGGT2015,KLN2006}; bounded derived length is the corresponding
obstruction in solvable families~\cite{LubotzkyWeiss1993}. Accordingly,
over prime fields outside a set of relative density zero, the uniform
expansion theorem of Breuillard--Gamburd~\cite{BreuillardGamburd2010}
improves the rate to $e^{-c4^m}$ (\Cref{cor:almost-all-primes}), and over
arbitrary finite fields the theorem of
Becker--Breuillard~\cite{BeckerBreuillard2026} bounds the number of
exceptional configurations at each fixed depth, up to simultaneous
conjugacy, by $O(q^\delta)$ for every $\delta>0$
(\Cref{cor:few-exceptional-blocks}).

\subsection*{Related work}
Word maps relate random arguments to the distribution of word
values~\cite{LarsenShalev,Shalev2013,BreuillardGlazer2025}. Every element
of a finite non-abelian simple group is a commutator~\cite{LOST2010};
iterated commutators at uniform arguments are asymptotically
equidistributed on simple groups of growing
order~\cite{GarionShalev2009}; the walk whose steps are uniform values of
a fixed word mixes in boundedly many steps~\cite{SchulShalev2011}; and
products of word values in disjoint variables are almost
uniform~\cite{LarsenShalevTiep2019}. These results concern a fixed word
at uniform arguments in simple groups of large order. Here the arguments
come from a prescribed conjugation-closed set, repeated occurrences are
read independently (\Cref{rem:notxy}), and the group is fixed while the
depth grows; the question is which subgroups obstruct contraction and at
what depth the obstruction disappears.

When $\mathcal C$ is a single class, the walk with steps uniform in
$\mathcal C$ is diagonalised by character ratios, and its mixing is
classical~\cite{DS81,Diaconis1988,Hildebrand1992,LulovPak,Roichman1996,
LiebeckShalev2005,Shalev2008}. In the commutator tree the steps are
prescribed, and trapping is detected through fixed vectors in induced
representations rather than through character bounds. In a different
direction, Amir, Blachar, Ghosh and Vishne~\cite{AmirMixability} ask
when some finite product of lazy steps, with freely chosen elements and
probabilities, is exactly uniform; every $S_n$ has this
property~\cite{AngelHolroyd2018}. We fix fair coins and the commutator
pattern and ask for contraction uniformly over valid configurations:
among the groups $S_n$, $n\ge5$, with transposition leaves, only $S_5$
has a level.

\subsection*{Origin}
The construction is Barrington's~\cite{Bar89}: width-five branching
programs over $S_5$ compute $\mathrm{NC}^1$ by combining their
instructions through a balanced tree of commutators, and over matrix
groups one may use pairs conjugate to their commutator~\cite{Boneh2013}. Group programs of
this kind are randomised in secure computation~\cite{Kilian1988}, and
the distribution of products of random group elements underlies
leakage-resilient circuits over $A_5$~\cite{MilesViola2013} and the
interleaved-product bounds of Gowers--Viola~\cite{GowersViola2019}. The
present question is what happens when each instruction of the expanded
program is independently replaced by the identity with probability
$1/2$. The companion paper~\cite{ZermenoS5} treats $S_5$ quantitatively;
here we establish the general structural law. The reduction from the
conditioned laws of branching programs to the walk studied here is
recalled in \Cref{sec:conditioning}.

\subsection*{Organisation}
\Cref{sec:prelim} sets up the model. \Cref{sec:criterion} proves the
fixed-vector criterion, \Cref{ithm:A} and the quantitative stability
estimates used later. \Cref{sec:synthesis} computes levels for symmetric
and alternating groups, minimal simple groups, transvections and
congruence quotients, and proves \Cref{icor:minsimple}.
\Cref{sec:rankone} proves \Cref{ithm:C}, \Cref{sec:octahedral-pair}
proves \Cref{ithm:B}, and \Cref{sec:quantitative-families} proves
\Cref{ithm:D}. An extended version of this paper~\cite{ZermenoExtended}
contains additional material that is not needed here: the localisation
of obstructions under quotients, the levels attained by prescribed
families of configurations, and further computational tables.
